\documentclass[11pt,a4paper]{article}
\usepackage[utf8]{inputenc}
\usepackage[T1]{fontenc}
\usepackage{amsmath,amssymb,amsthm}
\usepackage{booktabs,caption,microtype,xcolor}
\usepackage{pgfplots}\pgfplotsset{compat=1.18}
\usepackage[margin=2.5cm]{geometry}
\usepackage[colorlinks=true,linkcolor=blue!60!black,citecolor=blue!60!black,urlcolor=blue!60!black]{hyperref}
\captionsetup{font=small,labelfont=bf}
\newtheorem{theorem}{Theorem}
\newtheorem{proposition}[theorem]{Proposition}
\newtheorem{lemma}[theorem]{Lemma}
\newtheorem{conjecture}[theorem]{Conjecture}
\theoremstyle{remark}\newtheorem{remark}[theorem]{Remark}
\newcommand{\Z}{\mathbb{Z}}\newcommand{\Q}{\mathbb{Q}}\newcommand{\F}{\mathbb{F}}
\newcommand{\Off}{\mathrm{Off}}\newcommand{\Frob}{\operatorname{Frob}}
\newcommand{\target}[1]{\textcolor{red!70!black}{\small\sffamily[#1]}}

\title{The pair singular series of a polynomial III:\\ the off-diagonal in Ces\`aro form}
\author{Jasper Reichardt \and Claude Fable 5.1}
\date{Working draft, \today}

\begin{document}
\maketitle

\begin{abstract}
\target{Working draft. Theorem A is the target; nothing in Sections~3--4 is proved yet.}
Parts~I and~II reduced the conjecture $\sum_{h\le H}(S_f(h)-C(f)^2)=-\tfrac12C(f)\log H+A_f+o(1)$ for the pair
singular series of an irreducible polynomial $f$ to a statement about the off-diagonal sum over pairs of
distinct roots of $f$ modulo squarefree $d$ (``Hypothesis~(E)''), and showed that in Ces\`aro form any
unbounded saving over the trivial bound suffices when the Galois group is $2$-transitive. Here we prove
Hypothesis~(E) in Ces\`aro form for every irreducible quadratic, \target{goal}, by reducing it to weighted
Weyl sums of the roots of $\nu^2\equiv D\pmod d$ with a phase $e(kH/d)$, in the manner of Hooley's treatment of
$\sum_{n\le x}\tau(n^2+a)$; the frequency $k$ only needs to be controlled up to a power of $\log H$, because the
Ces\`aro weight makes the Fourier coefficients of the test function decay like $k^{-2}$. Over $\F_q[u]$ the
analogous statement, that the off-diagonal tends to zero as $q\to\infty$ at fixed degree, is a finite family
of point counts to which the Lang--Weil estimate applies. Numerically the off-diagonal of $t^2+1$ is bounded to
$10^7$ over $\Z$ and decreases with $q$ over $\F_q[u]$.
\end{abstract}

\section{The target and the reduction}\label{sec:target}

\paragraph{In plain terms.} The primes along a polynomial are correlated through the small primes; part~I
turned the strength of that correlation into an exact identity, whose main part (equal roots) is understood,
and whose remaining part (distinct roots modulo $d$) is a question about how the square roots of $D$ modulo
composite $d$ are distributed. That question has a sixty-year-old template, Hooley's divisor problem for
$n^2+a$; this paper adapts it. The weaker, Ces\`aro form of the statement is what the argument gives, and it is
enough for the leading term of the conjecture.

Notation is that of parts I and II: $f=at^2+bt+c$ irreducible with discriminant $D$, $\omega_f(p)$,
$W_f(d)=d\,b(d)\prod_{p\nmid d}P_p$ on squarefree $d$ with $|W_f(d)|\asymp\lambda(d)/d$, $\lambda$ multiplicative with
$\lambda(p)=p/(p-4)$ at split primes and $\lambda(p)=1$ at inert ones, $C=C(f)$, $\Sigma_f(H)=\sum_{h\le H}(S_f(h)-C^2)$.
Part~I, Theorem~6: with $\phi(\theta)=\theta-\theta^2$ and $m=m_{ss'}\in[1,d-1]$ the representative of $s'-s$,
\begin{equation}\label{eq:cesaro}
\sum_{h\le H}\Bigl(1-\frac hH\Bigr)\bigl(S_f(h)-C^2\bigr)=-\tfrac12C\log H+O_f(1)+\frac{C^2}H\Off^\ast_f(H),\qquad
\Off^\ast_f(H)=\sum_{d\ge2}W_f(d)\sum_{\{s,s'\}}\Bigl[-\frac{m(d-m)}d+\frac d2\Bigl(\phi\Bigl(\Bigl\{\frac{H-m}d\Bigr\}\Bigr)+\phi\Bigl(\Bigl\{\frac{H+m}d\Bigr\}\Bigr)\Bigr)\Bigr],
\end{equation}
the inner sum over unordered pairs of distinct roots modulo $d$, and trivially $\Off^\ast_f(H)\ll H\log H$.

\begin{theorem}[target]\label{thm:A}\target{TO PROVE}
For every irreducible quadratic $f$ there is $\eta>0$ with $\Off^\ast_f(H)\ll H(\log H)^{1-\eta}$. Consequently
\[
\sum_{h\le H}\Bigl(1-\frac hH\Bigr)\bigl(S_f(h)-C(f)^2\bigr)=-\tfrac12C(f)\log H+O\bigl((\log H)^{1-\eta}\bigr),
\]
i.e.\ Conjecture~1 of part~I holds in Ces\`aro form for every quadratic, with the first power of $C(f)$.
\end{theorem}

\subsection{Fourier expansion}

Since $\phi(\{x\})=\tfrac16-\pi^{-2}\sum_{k\ge1}k^{-2}\cos(2\pi kx)$ (absolutely convergent),
\begin{equation}\label{eq:fourier}
\Off^\ast_f(H)=\Omega_0-\frac1{\pi^2}\sum_{k\ge1}\frac1{k^2}\,T_k(H),\qquad
T_k(H)=\sum_{d}d\,W_f(d)\cos\Bigl(\frac{2\pi kH}d\Bigr)\sum_{\{s,s'\}}\cos\Bigl(\frac{2\pi km}d\Bigr),
\end{equation}
where $\Omega_0=\sum_dW_f(d)\sum_{\{s,s'\}}\bigl(\tfrac d6-\tfrac{m(d-m)}d\bigr)$ is independent of $H$. \target{check:
$\Omega_0$ converges absolutely (it is the second-moment discrepancy of the root differences, $\ll\sum_d\lambda(d)\,\omega(d)^2\,\mathrm{disc}_d$;
Kowalski--Soundararajan give $\mathrm{disc}_d\to0$ on average, but absolute convergence needs a rate; for
quadratics Hooley's Theorem~1 gives a power saving, which suffices).} The trivial bound $|T_k(H)|\ll H\log H$
comes from $d\,W_f(d)\asymp\lambda(d)$, so that with $K=(\log H)^{A}$ the frequencies $k>K$ contribute
$\ll H\log H/K$ to~\eqref{eq:fourier}, and Theorem~\ref{thm:A} follows from:

\begin{lemma}[A: twisted weighted Weyl sums]\label{lem:A}\target{TO PROVE}
There are $\eta,A>0$ such that for all $X,H\ge2$ and $1\le k\le(\log X)^{A}$,
\[
\sum_{d\le X}\lambda(d)\sum_{\nu^2\equiv D'\ (d)}e\Bigl(\frac{k(\nu+H)}d\Bigr)\ll X(\log X)^{1-\eta},\qquad D'=D/a^2,
\]
uniformly in $H$ (the phase $e(kH/d)$ is the point), together with the same bound for the moduli in
progressions $e\mid d$ with polynomial dependence on $e$, and the corresponding statement for the large
moduli $d>H$ after Hooley's reflection (Lemma~C below).
\end{lemma}

\emph{Why the frequency range is mild.} With the sawtooth $\{H/d\}$ of part~I's unsmoothed identity the
Fourier coefficients decay like $1/k$ and one would need $T_k$ to decay in $k$; the Ces\`aro weight replaces
$1/k$ by $1/k^2$, so it is enough to control $k\le(\log X)^{A}$ with any polynomial loss in $k$. This is the
single reason the Ces\`aro form is within reach and the sharp form is not.

\section{What is known about the pieces}\label{sec:known}

\paragraph{Unweighted, untwisted.} Hooley~\cite{Hooley1963}, Theorem~1: $\sum_{d\le X}\sum_{\nu^2\equiv D'(d)}e(k\nu/d)\ll_kX^{\theta}\log^2X$
with $\theta<1$, by parametrising the pairs $(\nu,d)$ through the primitive representations $d=Q(r,s)$ by a
finite set of binary quadratic forms of discriminant $4D'$ and reducing to incomplete Kloosterman sums (Weil's
bound). Hooley~\cite{Hooley1964} gives equidistribution without a power saving for any irreducible polynomial;
Kowalski--Soundararajan~\cite{KowalskiSoundararajan2021} reprove it as a mixing property of the Chinese
remainder theorem, with an average discrepancy $(\log X)^{-c}$.

\paragraph{Twisted, unweighted.} Hooley's Lemma~5 is a bound $R^\pm(h,X)=\sum_{k\le X}\rho(h,k)e(\pm hx/k)\ll\sigma(h)|h|X^{\theta}\log^3X$
for the same sums with the phase $e(\pm hx/k)$ built into the parametrisation \target{transcribe the exact
exponent from the original; the extracted text is illegible there}; it is exactly what his $\Sigma_5$, our
$\Off$, requires, and it is uniform in $x$.

\paragraph{Smooth weights, moduli in a progression.} Duke, Friedlander and Iwaniec~\cite{DFI2012}, Theorem~1.1:
for $f$ smooth on $[Y,2Y]$ with $|f|\le1$, $y^2|f''|\le1$, and either sign of $D$ (their \S16),
$\sum_{c\equiv0(q)}f(c)\sum_{b^2\equiv D(c)}e(hb/c)\ll h^{1/4}(Y+h\sqrt D)^{3/4}D^{1/8-1/1331}$, via the Kuznetsov
formula for half-integral-weight Kloosterman sums. Our phase $e(kH/c)$ has $c^2|\partial_c^2e(kH/c)|\asymp(kH/c)^2$, so
it satisfies their smoothness condition only for $Y\ge kH$: their theorem covers the moduli $c\ge kH$, not the
range $c\le H$ that carries the trivial bound. For $c\le H$ one needs Hooley's method with the phase, or a
version of Theorem~1.1 with $y^2|f''|\le\Delta^2$ and polynomial loss in $\Delta$ \target{check their proof for
the dependence on the derivative bound; the scaling $g(x)=\dots$ in their \S11 suggests the loss is
polynomial}.

\paragraph{Weights.} Write $\lambda=1*g$ with $g$ multiplicative, $g(p)=4/(p-4)$ at split primes and $0$
otherwise; then $\sum_{d\le X}\lambda(d)F(d)=\sum_{e}g(e)\sum_{m\le X/e}F(em)$, the moduli $e>E$ contribute
$\ll E^{-1+\varepsilon}X\log X$, and only $e\le(\log X)^{B}$ matter. So Lemma~A reduces to the unweighted sums over
moduli $\equiv0\bmod e$ with polynomial dependence on $e$ (Iwaniec~\cite{Iwaniec1978} for $D=-1$; Ngo for $D>0$).

\paragraph{Large moduli.} For $d>H$ a pair of distinct roots modulo $d$ has $m\le H$ only if every prime
factor of $d$ divides $hQ(h)$ for some $h\le H$, $Q(h)=a^2h^2-D$; Hooley's reflection writes $d=Q(h)/e$ with
$e<Q(h)/H$, turning the large moduli into a sum over the small divisors $e$ of $Q(h)$ (part~I, companion note,
Prop.~A3), which is the same structure with $e$ in place of $d$ and the level-of-distribution input for
$\#\{h\le H:e\mid Q(h)\}$ from~\cite{Iwaniec1978} and Grimmelt--Merikoski.

\section{Plan of the proof of Lemma~A}\label{sec:plan}
\target{To be written. Steps: (1) Dirichlet convolution to remove $\lambda$; (2) dyadic decomposition in $d$;
(3) $d\ge kH$: DFI Theorem~1.1 with the phase absorbed in $f$; (4) $d<kH$: Hooley's parametrisation
$\nu^2-D'=d\ell$, $d=Q(r,s)$ primitive, $\nu/d\equiv\bar s\,(\dots)/Q(r,s)$; the phase $e(kH/Q(r,s))$ is slowly
varying in $(r,s)$ over the regions where Hooley bounds the incomplete Kloosterman sums, so it can be removed by
partial summation at the cost of a factor $kH/Y$ per region unless $Y\ge(kH)^{1/2}$, and below that a genuinely
new estimate is needed: this is the open core. (5) Assemble with $K=(\log H)^{A}$.}

\section{Route B: the same statement over $\F_q[u]$}\label{sec:B}

Let $\mathcal A=\F_q[u]$, $q$ odd, $f=t^2-D$ with $D\in\mathcal A$ squarefree, and for $N\ge1$ let $\Off_f(N)$ be the
off-diagonal of the identity of part~I, Theorem~8 (the sum over the $q^N-1$ nonzero $h$ of degree $<N$). Part~I
proved that the moduli of degree $\le N$ contribute nothing and that the remainder is a finite sum:
\[
\Off_f(N)=\sum_{N<\deg d\le M(N)}W_f(d)\,A_N(d)-q^N\sum_{\deg d>N}\frac{W_f(d)(\omega_f(d)^2-\omega_f(d))}{|d|},\qquad M(N)=N-1+\max(2N-2,\deg D),
\]
$A_N(d)$ the number of ordered pairs of distinct roots modulo $d$ whose difference has degree $<N$.

\begin{theorem}[target]\label{thm:B}\target{TO PROVE}
For fixed $N$ and $\deg D$, $\Off_f(N)=O(q^{-1/2})$ as $q\to\infty$. Hence the function-field form of
Conjecture~1, $\sum_{h}(S_f(h)-C^2)=-C(f)N+O(1)$, holds with $O(1)$ replaced by $O(q^{-1/2})$.
\end{theorem}

\emph{Reduction to point counts.} For each degree $k\in(N,M(N)]$, the sum $\sum_{\deg d=k}W_f(d)A_N(d)$ is a
weighted count of triples $(d,s,s')$ with $d$ squarefree monic of degree $k$, $s^2\equiv s'^2\equiv D\pmod d$,
$s\ne s'$, and $\deg(\widetilde{s'-s})<N$. Writing $s'-s=m$ with $d\mid m^2-4D$ on the part of $d$ where the roots
differ and $d\mid m$ on the rest (part~I, Lemma~9), the condition is algebraic in the coefficients of $d$ and
$m$, so the count is $\#V_{k,N}(\F_q)$ for an explicit variety $V_{k,N}$ defined over $\F_q$ (depending on $D$),
weighted by $W_f(d)=P(1)\lambda(d)/|d|$, and $\lambda(d)$ is a product of local factors $1+O(q^{-\deg P})$ over the
prime factors of $d$. The expected term $q^N\sum_{\deg d>N}W_f(d)\omega^\circ_f(d)/|d|$ is the same count with the
probability $q^{N-k}$ in place of the indicator of $\deg m<N$. Theorem~\ref{thm:B} therefore follows from the
Lang--Weil estimate once one knows, for each of the finitely many $(k,N)$, that the components of $V_{k,N}$ over
$\overline{\F}_q$ have the dimensions and multiplicities that make the main term equal to the expected one;
this is a finite algebraic computation for each $N$, and a uniform statement is the content of Katz-type
equidistribution for the Sali\'e sums of the roots over primes of fixed degree.

\IfFileExists{data/offq.dat}{%
\begin{figure}[htbp]\centering
\begin{tikzpicture}
\begin{axis}[width=0.7\textwidth,height=5.2cm,xlabel={$q$},ylabel={$\Off_f(N)$},legend pos=south east,legend style={font=\scriptsize},
  title={$f=t^2-u$ over $\F_q[u]$: the off-diagonal at fixed $N$ as $q$ grows}]
\addplot[only marks,mark=*,blue!70!black] table[x=q,y=off,row sep=newline] {data/offq-N2.dat};
\addplot[only marks,mark=square*,red!70!black] table[x=q,y=off] {data/offq-N3.dat};
\addplot[dashed,gray] coordinates {(3,0) (47,0)};
\legend{$N=2$,$N=3$}
\end{axis}
\end{tikzpicture}
\caption{Exact values of $\Off_f(N)$ over $\F_q[u]$ for $f=t^2-u$ (\texttt{scripts/offq.ts}). Theorem~\ref{thm:B}
predicts decay like $q^{-1/2}$.}
\label{fig:offq}
\end{figure}}{}

\section{Numerical evidence}\label{sec:num}
Over $\Z$: the off-diagonal of $t^2+1$ is bounded to $H=10^7$, with mean $-0.01$, standard deviation $0.51$ and
no drift in $\log H$ (part~II, Section~5.3); a failure of Hypothesis~(E) would show as a slope of order $0.73$
per unit of $\log H$. Over $\F_q[u]$: Figure~\ref{fig:offq}.

\section{Open issues}
(1)~The core of Lemma~A: the phase $e(kH/d)$ for $d<(kH)^{1/2}$ within Hooley's parametrisation. (2)~The
absolute convergence of $\Omega_0$ with a rate. (3)~Theorem~B: identify the components of $V_{k,N}$ for $N=2,3$
and $\deg D=1$ by hand, then the general pattern. (4)~Non-$2$-transitive Galois groups, where a saving over the
trivial bound is not enough (part~I, Remark~7).

\bibliographystyle{plain}
\bibliography{refs}
\end{document}
